\section{The fan of a belt cell}
\label{sec:area_law_cell_fans}

The local fan is constructed before its triangles are colored. Each side
may remain whole or be divided at its midpoint, according to the corners
of the opposing cells. The following construction permits every such
choice of subdivisions; see \cite[Section~11]{OpenAI2026AreaLaw}.

% Source: 10-geometry.tex, lines 299--310. This is the single-cell
% geometric construction. Opposing-region assignments and merged
% identifiers are subsequent constructions.

\begin{definition}[Elementary segments and fan triangles]
    \label{def:area_law_cell_fan}
    \lean{TNLean.PEPS.AreaLaw.Geometry.CellFanSlot,
        TNLean.PEPS.AreaLaw.Geometry.cellFanCenter,
        TNLean.PEPS.AreaLaw.Geometry.cellFanStart,
        TNLean.PEPS.AreaLaw.Geometry.cellFanEnd,
        TNLean.PEPS.AreaLaw.Geometry.cellFanPolygon}
    \leanok
    \uses{def:area_law_dyadic_cells,def:area_law_template}
    Let $Q$ be a closed dyadic square of side $t=2^\ell$ and center $c$.
    On each of its four sides choose either the whole side or its two
    half sides. The resulting segments are called elementary segments.
    For each elementary segment $e=[a_e,b_e]$, let
    \begin{align}
        P_e=\operatorname{conv}\{c,a_e,b_e\}.
    \end{align}
    These are nondegenerate closed triangles. Their perimeter and radial
    edges have slopes $0,\infty,1,-1$.
\end{definition}

\begin{theorem}[A dyadic square as a maximum-norm ball]
    \label{thm:area_law_cell_fan_closed_ball}
    \lean{TNLean.PEPS.AreaLaw.Geometry.closedBall_cellFanCenter_eq_closure_dyadicCell}
    \leanok
    \uses{def:area_law_dyadic_cells,def:area_law_cell_fan}
    Let $o\in\R^2$, $\ell\in\N$ and $z\in\Z^2$. Let
    $Q=o+2^\ell(z+[0,1)^2)$ be the half-open dyadic cell, with
    center $c=o+2^\ell(z+(1/2,1/2))$ and half-side $r=2^\ell/2$.
    Then
    \begin{align}
        \overline{B}_\infty(c,r) & =\overline Q.
    \end{align}
\end{theorem}
\begin{proof}\leanok
    \uses{thm:area_law_dyadic_closures}
    The coordinate description of $\overline Q$ gives
    $x\in\overline Q$ if and only if $|x_j-c_j|\le r$ for $j=1,2$.
    These inequalities are equivalent to $\|x-c\|_\infty\le r$.
\end{proof}

\begin{theorem}[Fan vertices are actual cell marks]
    \label{thm:area_law_cell_fan_mark_vertices}
    \lean{TNLean.PEPS.AreaLaw.Geometry.cellFan_vertices_mem_beltCellMarks}
    \leanok
    \uses{def:area_law_cell_fan,def:area_law_belt_cell_marks}
    The center and both endpoints of every elementary segment belong
    to the actual nine-point mark set $M_{\ell,z}$ of the same cell.
\end{theorem}
\begin{proof}\leanok
    \uses{def:area_law_cell_fan,def:area_law_belt_cell_marks}
    The center has relative coordinates $(1/2,1/2)$. Every segment
    endpoint is a corner or side midpoint, so both relative coordinates
    are among $0,1/2,1$.
\end{proof}

\begin{theorem}[Number of fan triangles]
    \label{thm:area_law_cell_fan_count}
    \lean{TNLean.PEPS.AreaLaw.Geometry.card_cellFanSlot_bounds}
    \leanok
    \uses{def:area_law_cell_fan}
    Every such fan has between four and eight triangles.
\end{theorem}
\begin{proof}\leanok
    \uses{def:area_law_cell_fan}
    Each of the four sides contributes either one or two elementary
    segments.
\end{proof}

\begin{theorem}[Exact cover by fan triangles]
    \label{thm:area_law_cell_fan_cover}
    \lean{TNLean.PEPS.AreaLaw.Geometry.cellFanPolygons_cover}
    \leanok
    \uses{def:area_law_cell_fan}
    The closed fan triangles cover exactly the closed dyadic square:
    \begin{align}
        Q=\bigcup_e P_e.
    \end{align}
\end{theorem}
\begin{proof}\leanok
    \uses{def:area_law_cell_fan}
    Every point of $Q$ lies on the segment from $c$ to some point of the
    perimeter. Every perimeter point belongs to an elementary segment,
    so it is in the corresponding fan triangle. Conversely, every
    vertex lies in $Q$, and convexity keeps the triangle in $Q$.
\end{proof}

\begin{theorem}[Intersections of fan triangles]
    \label{thm:area_law_cell_fan_intersections}
    \lean{TNLean.PEPS.AreaLaw.Geometry.cellFanPolygons_inter_eq}
    \leanok
    \uses{def:area_law_cell_fan}
    For any two elementary segments $e,e'$ of the same fan,
    \begin{align}
        P_e\cap P_{e'}
        =\{c\}\cup\bigcup_{u\in e\cap e'}[c,u].
    \end{align}
    Thus a shared perimeter endpoint gives a shared radial segment,
    while disjoint elementary segments give intersection $\{c\}$.
\end{theorem}
\begin{proof}\leanok
    \uses{def:area_law_cell_fan}
    A point different from $c$ determines its unique perimeter endpoint
    on the outward ray from $c$. It belongs to both triangles exactly
    when that endpoint lies in both elementary segments. The center
    belongs to every triangle.
\end{proof}


\subsection{The primary containing a non-belt fine cell}

\begin{definition}[The shifted pitch index]
    \label{def:area_law_nonbelt_pitch_index}
    \lean{TNLean.PEPS.AreaLaw.Geometry.nonbeltPitchIndex}
    \leanok
    \uses{def:area_law_primary_regions}
    For nonnegative fine and pitch exponents $\ell\le p$, integer shift representatives
    $a,b$ and fine-cell index $z$, put $m=2^{p-\ell}$ and define
    \begin{align}
        j(z)=\left(\left\lfloor\frac{z_1-a}{m}\right\rfloor,
        \left\lfloor\frac{z_2-b}{m}\right\rfloor\right).
    \end{align}
    These are the coordinate indices of the shifted pitch square.
\end{definition}

\begin{theorem}[Closed-cell containment in an actual primary]
    \label{thm:area_law_nonbelt_primary_containment}
    \lean{TNLean.PEPS.AreaLaw.Geometry.closure_nonbeltCell_subset_primaryBirthRegion}
    \leanok
    \uses{def:area_law_nonbelt_pitch_index,def:area_law_primary_regions,
        def:area_law_fine_layer_indices,def:area_law_belt_cells}
    Fix an origin $o$, finite endpoint set $Z$, layer index $k$ and
    nonnegative integer width $C_0$. Suppose that $\ell\le k$ and
    $\ell\le p$. Choose residues $0\le a,b<m=2^{p-\ell}$, and let $z$
    index an actual fine cell in this layer that is omitted by both
    selected belt residues. Its closed square satisfies
    \begin{align}
        \overline{Q_\ell(z)}\subseteq
        \overline{D_k\cap I_{a,b}(j(z))},
    \end{align}
    where $I_{a,b}(j)$ is the corresponding open pitch interior.
\end{theorem}
\begin{proof}\leanok
    \uses{def:area_law_nonbelt_pitch_index,def:area_law_primary_regions,
        thm:area_law_fine_cell_containment,thm:area_law_dyadic_closures}
    Since neither residue is selected, each shifted coordinate remainder
    lies between $1$ and $m-1$. The open fine-cell interior consequently
    lies strictly inside the indicated pitch interior and belongs to the
    actual layer. Its closure is the whole closed fine square. Taking
    closures gives the inclusion, including when a fine-cell side
    touches a pitch-interior boundary.
\end{proof}

\begin{theorem}[Unique actual primary identifier]
    \label{thm:area_law_nonbelt_unique_primary}
    \lean{TNLean.PEPS.AreaLaw.Geometry.exists_unique_primary_of_nonbeltCell}
    \leanok
    \uses{thm:area_law_nonbelt_primary_containment,
        def:area_law_primary_pitch_indices}
    Under the preceding hypotheses, there is exactly one retained
    primary pitch index whose birth region contains the whole closed
    fine square. This index is $j(z)$.
\end{theorem}
\begin{proof}\leanok
    \uses{thm:area_law_nonbelt_primary_containment,
        thm:area_law_primary_pitch_membership,thm:area_law_primary_diameter_separation}
    The open fine square is nonempty and lies in the actual
    layer--pitch intersection, so the pitch index is retained.
    If two different primary birth regions contained this square,
    an interior point would belong to both. Their positive same-layer
    separation rules this out.
\end{proof}

% Source: 10-geometry.tex, prop:two-families, lines 299--323,
% especially 308--316, and geometry:initial-stars, lines 352--363;
% revision adc7f1241b42e322a6451854ab7e4b4c146bf78a.

\begin{theorem}[Constant radius of an elementary base]
    \label{thm:area_law_cell_fan_base_radius}
    \lean{TNLean.PEPS.AreaLaw.Geometry.norm_sub_cellFanCenter_of_mem_base}
    \leanok
    \uses{def:area_law_cell_fan}
    For any origin, natural exponent $\ell$, integer cell index and
    optional midpoint subdivisions, let $c$ be the cell center and
    $e=[a_e,b_e]$ an actual elementary base. Every $x\in e$ satisfies
    \begin{align}
        \|x-c\|_\infty & =\frac{2^\ell}{2}.
    \end{align}
    This is the actual base-radius identity in the fan construction
    of \cite[Section~11]{OpenAI2026AreaLaw}.
\end{theorem}
\begin{proof}\leanok
    \uses{def:area_law_cell_fan}
    In the affine parametrization of a side, one relative coordinate
    is $\pm2^\ell/2$ and the other is $(2^\ell/2)u$ with $|u|\le1$.
    The whole side and either midpoint half all satisfy this bound.
    The maximum of the two absolute coordinates is therefore $2^\ell/2$.
\end{proof}
